%% =============================================================
%% Red Hat Technical Document Template
%% Engine: pdflatex (Overleaf default) -- NO xelatex required
%%
%% Usage on Overleaf:
%%   1. Upload this .tex file to a new Overleaf project
%%   2. Leave compiler as pdflatex (default)
%%   3. Click Compile
%% =============================================================
\documentclass[11pt,a4paper]{article}

%% ---- Encoding & Language ----
\usepackage[utf8]{inputenc}
\usepackage[T1]{fontenc}
\usepackage[english]{babel}

%% ---- pdflatex-compatible Fonts ----
%% Helvetica clone (geometric sans-serif, close to Red Hat Display)
\usepackage[scaled=0.92]{helvet}
\usepackage{courier}
\renewcommand{\familydefault}{\sfdefault}

%% ---- Page Layout ----
\usepackage[a4paper,
            left=25mm, right=25mm,
            top=30mm,  bottom=30mm,
            headheight=20mm]{geometry}

%% ---- Colors ----
\usepackage{xcolor}

%% ---- Graphics ----
\usepackage{graphicx}

%% ---- Section Titles ----
\usepackage{titlesec}

%% ---- Headers & Footers ----
\usepackage{fancyhdr}

%% ---- Hyperlinks ----
\usepackage{hyperref}

%% ---- Callout Boxes (explicit libraries only -- avoids fontspec on TeX Live 2025) ----
\usepackage{tcolorbox}
\tcbuselibrary{skins, breakable}

%% ---- Code Listings ----
\usepackage{listings}

%% ---- Tables ----
\usepackage{tabularx}
\usepackage{booktabs}

%% ---- Lists ----
\usepackage{enumitem}

%% ---- Placeholder Text ----
\usepackage{lipsum}

%% ---- Title Page Background Color ----
\usepackage{pagecolor}
\usepackage{afterpage}

%% =============================================================
%%  RED HAT BRAND COLORS
%% =============================================================
\definecolor{RHRed}{HTML}{EE0000}
\definecolor{RHBlack}{HTML}{000000}
\definecolor{RHWhite}{HTML}{FFFFFF}
\definecolor{RHDarkGray}{HTML}{292929}
\definecolor{RHLightGray}{HTML}{E0E0E0}
\definecolor{RHDarkPurple}{HTML}{21134D}
\definecolor{RHCodeBg}{HTML}{1E1E1E}
\definecolor{RHCodeFg}{HTML}{D4D4D4}
\definecolor{RHGreen}{HTML}{3E8914}
\definecolor{RHCyan}{HTML}{007A87}

%% =============================================================
%%  SECTION TITLE FORMATTING
%% =============================================================
\titleformat{\section}
  {\Large\bfseries\color{RHRed}}{\thesection}{1em}{}
  [\color{RHRed}\titlerule]
\titleformat{\subsection}
  {\large\bfseries\color{RHBlack}}{\thesubsection}{1em}{}
\titleformat{\subsubsection}
  {\normalsize\bfseries\color{RHDarkGray}}{\thesubsubsection}{1em}{}

%% =============================================================
%%  HEADER AND FOOTER
%% =============================================================
\pagestyle{fancy}
\fancyhf{}
\fancyhead[L]{\small\textbf{\textcolor{RHDarkGray}{Red Hat Technical Document}}}
\fancyhead[R]{\small\textbf{\textcolor{RHRed}{PUBLIC}}}
\fancyfoot[L]{\small\textcolor{RHDarkGray}{\copyright~\the\year~Red Hat, Inc.}}
\fancyfoot[R]{\small\textcolor{RHDarkGray}{\thepage}}
\renewcommand{\headrulewidth}{1.5pt}
\renewcommand{\headrule}{%
  \hbox to\headwidth{\color{RHRed}\leaders\hrule height \headrulewidth\hfill}}
\renewcommand{\footrulewidth}{0.4pt}
\renewcommand{\footrule}{%
  \hbox to\headwidth{\color{RHLightGray}\leaders\hrule height \footrulewidth\hfill}}

%% =============================================================
%%  HYPERLINKS
%% =============================================================
\hypersetup{
    colorlinks = true,
    linkcolor  = RHRed,
    urlcolor   = RHRed,
    citecolor  = RHRed,
    pdftitle   = {Red Hat Technical Document Template},
    pdfauthor  = {Gabriel Tavares},
    pdfsubject = {Red Hat LaTeX Template -- Overleaf Compatible}
}

%% =============================================================
%%  CUSTOM CALLOUT BOXES
%%  Using only 'skins' and 'breakable' -- safe on all TeX Live versions
%% =============================================================
\newtcolorbox{rhnote}{
    enhanced, breakable,
    colback   = RHLightGray!25,
    colframe  = RHRed,
    coltitle  = RHWhite,
    fonttitle = \bfseries,
    title     = {Note},
    attach boxed title to top left = {yshift=-2mm, xshift=4mm},
    boxed title style = {colback=RHRed, colframe=RHRed, sharp corners}
}

\newtcolorbox{rhwarning}{
    enhanced, breakable,
    colback   = RHRed!6,
    colframe  = RHRed,
    coltitle  = RHWhite,
    fonttitle = \bfseries,
    title     = {Warning},
    attach boxed title to top left = {yshift=-2mm, xshift=4mm},
    boxed title style = {colback=RHRed, colframe=RHRed, sharp corners}
}

\newtcolorbox{rhtip}{
    enhanced, breakable,
    colback   = RHDarkPurple!6,
    colframe  = RHDarkPurple,
    coltitle  = RHWhite,
    fonttitle = \bfseries,
    title     = {Tip},
    attach boxed title to top left = {yshift=-2mm, xshift=4mm},
    boxed title style = {colback=RHDarkPurple, colframe=RHDarkPurple, sharp corners}
}

%% =============================================================
%%  CODE LISTINGS  (dark theme, no minted)
%% =============================================================
\lstset{
    backgroundcolor  = \color{RHCodeBg},
    basicstyle       = \ttfamily\small\color{RHCodeFg},
    keywordstyle     = \color{RHCyan}\bfseries,
    commentstyle     = \color{RHGreen},
    stringstyle      = \color{orange},
    frame            = single,
    rulecolor        = \color{RHRed},
    breaklines       = true,
    showstringspaces = false,
    numbers          = left,
    numberstyle      = \tiny\color{RHDarkGray},
    numbersep        = 8pt,
    xleftmargin      = 12pt
}

%% ---- Bullet style: red square marker ----
\setlist[itemize]{label=\textcolor{RHRed}{\rule{0.6em}{0.6em}}}

%% =============================================================
%%  DOCUMENT BEGINS
%% =============================================================
\begin{document}

%% ---- Title Page ----
\newpagecolor{RHRed}
\thispagestyle{empty}
\color{RHWhite}

\vspace*{2.5cm}

{\fontsize{28}{34}\selectfont\bfseries Red Hat Technical\\[0.3cm]Document Template}

\vspace{0.6cm}
\noindent\textcolor{RHWhite}{\rule{\textwidth}{1pt}}
\vspace{0.4cm}

{\large A Professional \LaTeX{} Template\\following Red Hat Brand Guidelines}

\vfill

\begin{tabular}{ll}
    \textbf{Author:}  & Gabriel Tavares \\[2pt]
    \textbf{Company:} & Red Hat, Inc.   \\[2pt]
    \textbf{Date:}    & \today          \\[2pt]
    \textbf{Version:} & 1.0             \\
\end{tabular}

\vspace{1.5cm}
\begin{flushright}
    {\LARGE\bfseries RED HAT}
\end{flushright}

\afterpage{\restorepagecolor}
\newpage
\color{RHBlack}

%% ---- Table of Contents ----
\tableofcontents
\newpage

%% =============================================================
\section{Introduction}
%% =============================================================

This document serves as a professional technical writing template aligned with Red Hat's corporate brand guidelines. It is designed to compile with \textbf{pdflatex}, making it fully compatible with \href{https://www.overleaf.com}{Overleaf} using the default engine without any configuration changes.

The template provides a consistent visual identity for whitepapers, architecture guides, solution briefs, and internal technical reports. All color definitions, typography choices, and layout conventions follow the official Red Hat Brand Standards.

\subsection{Brand Colors}

The following colors are pre-configured and available throughout the document:

\begin{center}
\begin{tabularx}{\textwidth}{llX}
    \toprule
    \textbf{Color Name} & \textbf{Hex Code} & \textbf{Usage} \\
    \midrule
    \textcolor{RHRed}{\rule{1em}{1em}}~Red Hat Red    & \texttt{\#EE0000} & Primary accent, section headers, borders \\
    \textcolor{RHBlack}{\rule{1em}{1em}}~Black         & \texttt{\#000000} & Body text, headings \\
    \textcolor{RHDarkGray}{\rule{1em}{1em}}~Dark Gray  & \texttt{\#292929} & Secondary text, footers \\
    \textcolor{RHLightGray}{\rule{1em}{1em}}~Light Gray & \texttt{\#E0E0E0} & Backgrounds, callout boxes \\
    \textcolor{RHDarkPurple}{\rule{1em}{1em}}~Dark Purple & \texttt{\#21134D} & Alternative accent, tip boxes \\
    \bottomrule
\end{tabularx}
\end{center}

%% =============================================================
\section{Callout Boxes}
%% =============================================================

Three custom callout box environments are available to draw the reader's attention to important information.

\begin{rhnote}
This is a \textbf{Note} box. Use it to provide supplementary information, clarifications, or context that is helpful but not critical to the main flow of the document.
\end{rhnote}

\vspace{0.5em}

\begin{rhwarning}
This is a \textbf{Warning} box. Use it to alert the reader to potential issues, risks, or actions that could have negative consequences if not followed carefully.
\end{rhwarning}

\vspace{0.5em}

\begin{rhtip}
This is a \textbf{Tip} box. Use it to share best practices, shortcuts, or recommended approaches that can improve the reader's workflow or outcome.
\end{rhtip}

%% =============================================================
\section{Code Snippets}
%% =============================================================

Technical documents frequently require the inclusion of code, configuration files, or command-line instructions. The \texttt{lstlisting} environment is pre-configured with a dark theme consistent with Red Hat's developer tooling aesthetics.

\begin{lstlisting}[language=bash, caption={Deploying an application on OpenShift}]
# Login to OpenShift cluster
oc login --token=<your-token> \
         --server=https://api.cluster.example.com:6443

# Create a new project
oc new-project my-application

# Deploy from a container image
oc new-app \
  --image=registry.access.redhat.com/ubi9/ubi:latest

# Expose the service
oc expose svc/ubi
\end{lstlisting}

\begin{lstlisting}[language=bash, caption={Ansible Playbook -- Install and configure HTTPD}]
---
- name: Install and configure HTTPD
  hosts: webservers
  become: true
  tasks:
    - name: Install HTTPD package
      ansible.builtin.yum:
        name: httpd
        state: present

    - name: Start and enable HTTPD service
      ansible.builtin.service:
        name: httpd
        state: started
        enabled: true
\end{lstlisting}

%% =============================================================
\section{Technical Architecture}
%% =============================================================

\subsection{Overview}

\lipsum[1]

\subsection{Component Breakdown}

The following table provides a summary of the key components in a typical Red Hat hybrid cloud architecture:

\begin{center}
\begin{tabularx}{\textwidth}{llX}
    \toprule
    \textbf{Component} & \textbf{Product} & \textbf{Description} \\
    \midrule
    Container Platform & OpenShift & Enterprise Kubernetes platform for application deployment and management \\
    Automation         & Ansible   & Agentless IT automation for provisioning, configuration, and orchestration \\
    Operating System   & RHEL      & Enterprise Linux providing a stable, secure, and certified foundation \\
    Service Mesh       & OpenShift Service Mesh & Traffic management and observability for microservices \\
    \bottomrule
\end{tabularx}
\end{center}

\subsection{Key Considerations}

When designing a solution based on Red Hat technologies, the following principles should guide architectural decisions:

\begin{itemize}
    \item \textbf{Security by Default:} All Red Hat products are built with security in mind, leveraging SELinux, FIPS compliance, and supply chain security.
    \item \textbf{Open Standards:} Red Hat's commitment to open source ensures interoperability and avoids vendor lock-in.
    \item \textbf{Hybrid Cloud Portability:} Applications should be designed to run consistently across on-premises, public cloud, and edge environments.
    \item \textbf{Automation First:} Leverage Ansible Automation Platform to reduce manual toil and enforce consistent configurations at scale.
\end{itemize}

%% =============================================================
\section{Conclusion}
%% =============================================================

\lipsum[3]

\vspace{1.5em}
\noindent\textcolor{RHRed}{\rule{\textwidth}{1.5pt}}

\vspace{0.5em}
\begin{center}
    \large\textbf{\textcolor{RHRed}{RED HAT}} \\[2pt]
    \small\textcolor{RHDarkGray}{Be open. Be brave. Be Red Hat.}
\end{center}

\end{document}
